% EES Latex-Vorlage Januar 2021

\documentclass[
12pt,               										% Schriftgr��e
a4paper,            										% Layout f�r DINA4
german,             										% deutsche Sprache, global
twoside,            										% Layout f�r beidseitigen Druck
headinclude,        									% Kopfzeile wird Seiten-Layouts mit ber�cksichtigt
headsepline,        									% horizontale Linie unter Kolumnentitel
plainheadsepline,										% horizontale Linie unter Kolumnentitel auch bei Chapter
BCOR=20mm,           									% Korrektur f�r die Bindung
DIV=18,              									% DIV-Wert f�r die Erstellung des Satzspiegels, siehe scrguide
parskip=half,       									% Absatzabstand statt Absatzeinzug
openany,            										% Kapitel k�nnen auf geraden und ungeraden Seiten beginnen
bibliography=totoc,version=first, 						% Literaturverz. wird ins Inhaltsverzeichnis eingetragen
numbers=noenddot,   								% Kapitelnummern immer ohne Punkt
captions=tableheading,version=first, 					% korrekte Abst�nde bei Tabellen�berschriften
fleqn,             										% fleqn: Glgen links (statt mittig)
listof=totoc,version=first,								% Abbildungs- und Tabellenverzeichnis ins Inhaltsverzeichnis
cleardoublepage=empty								% Kopfzeile bei leeren Seiten entfernen
]{scrbook}


%--------------- Packages ----------------
\usepackage[ngerman]{babel}							% Neue deutsche Trennmuster
\usepackage[utf8]{inputenc}							% direkte Eingabe von Umlauten & Co.
\usepackage[T1]{fontenc}								% T1-Schriften
\usepackage{lmodern}								% moderne Schriften
\usepackage[format=hang,							% Captions ausrichten
justification=raggedright,
singlelinecheck=off,
labelfont=bf,
font=small,
skip=4pt
]{caption} 											
\usepackage[centertags]{amsmath} 					% AMS-Mathematik, centertags zentriert Nummer bei split
\usepackage{tabularx}								% erweiterte Tabellen
\usepackage{graphicx}            							% zum Einbinden von Grafiken
\usepackage{setspace}            							% Zeilenabstand einstellbar
\usepackage{scrlayer-scrpage}           					% Kopf- und Fu�zeilen-Layout 
\usepackage[pdfborderstyle={/S/U/W 1}, ]{hyperref}		% Die Links im PDF werden nur noch d�nn und rot unterstrichen, nicht mehr fett umrahmt. 
\usepackage[absolute]{textpos}						% benutzerdefinierte Textpositionierung 
\usepackage{charter}

\hypersetup{%
	colorlinks=true, 		% Aktivieren von farbigen Links im Dokument
	linkcolor=black, 	% Farbe festlegen
	citecolor=black,
	filecolor=black,
	menucolor=black,
	urlcolor=black,
	linktocpage=true, 		% Nicht der Text sondern die Seitenzahlen in Verzeichnissen klickbar
	bookmarksnumbered=true 	% Überschriftsnummerierung im PDF Inhalt anzeigen.
}

%--------------- Sonstiges ----------------

\pagestyle{scrheadings}									% Kopf- und Fu�zeile...
\renewcommand{\headfont}{\normalfont\sffamily}    		% Kolumnentitel serifenlos
\renewcommand{\pnumfont}{\normalfont\sffamily}    		% Seitennummern serifenlos
\ihead[]{\headmark}              								% Kopfzeile innen
\ohead[\pagemark]{\pagemark}     							% Kopfzeile au�en
\ifoot[]{} 												% Fu�zeile innen
\ofoot[]{}    												% Fu�zeile au�en	
\setlength{\headheight}{1.5\baselineskip}
\onehalfspacing											% 1,5 Zeilenabstand
%\typearea[current]{current}        							% Neuberechnung des Satzspiegels mit alten Werten nach �nderung von Zeilenabstand,etc
\renewcommand{\bibname}{Literatur und Quellen} 			% Literaturverzeichnisbezeichnung
\renewcommand{\figurename}{Abb.}   						% Abbildungsbezeichnung
\renewcommand{\listfigurename}{Abbildungsverzeichnis} 	% Abbildungsverzeichnisbezeichnung
\renewcommand{\captionfont}{\small}						% Bildunterschriften klein kursiv
\graphicspath{{figs/}{bilder/}}    							% Bildverzeichnis

\RedeclareSectionCommand[beforeskip=0pt]{chapter}		%kein Abstand zum oberen Seitenrand bei neuen Kapiteln


\makeindex


%----------------------- Beginn des Dokuments -----------------------
\begin{document}
\pagenumbering{Roman}									% R�mische Seitennummerierung

\thispagestyle{empty}
\setcounter{page}{-1}

\begin{textblock*}{\textwidth}(30mm,33mm)
\includegraphics[width=47mm]{ees}
\end{textblock*}

\begin{textblock*}{\textwidth}(127mm,192mm)
\includegraphics[width=93mm]{fausiegel}
\end{textblock*}

{\raggedleft
\textsc{Lehrstuhl für Elektrische Energiesysteme}\\
Vorstand: Univ.-Prof. Dr.-Ing. Matthias Luther
\par}

\vspace{51mm}

{\centering

%##########################################################
% BEARBEITUNGSFELDER
\large{Seminararbeit} \\	
\Large{Lastflussberechnung von Umrichter-dominierten Netzen}
%##########################################################

\par}

\vspace{105mm}

{\raggedright
\begin{tabbing}
XX \= XXXXXXXXXXXXX \= XXXXXXXXXXXXXXXXXXXXXX \kill
%##########################################################
% BEARBEITUNGSFELDER
		\> \textbf{Bearbeiter:} 	\> B. Eng. Maximilian Köhler \\
		\>						\> 23176975		\\
 		\>												\>									\\
		\> \textbf{Betreuer:}		\> M. Sc. Dominik Frauenknecht \\
		\>												\>									\\
		\> \textbf{Abgabedatum:}	\> 31.03.2024															
%##########################################################
\end{tabbing}
\par}
\cleardoublepage									% Deckblatt
Ich versichere, dass ich die vorliegende Arbeit ohne fremde Hilfe und ohne Benutzung anderer als der angegebenen Quellen angefertigt habe, und dass die Arbeit in gleicher oder ähnlicher Form noch keiner anderen Prüfungsbehörde vorgelegen hat und von dieser als Teil einer Prüfungsleistung angenommen wurde. Alle Ausführungen, die wörtlich oder sinngemäß übernommen wurden, sind als solche gekennzeichnet.
\vskip 0.5cm

\begin{table}[!htb]
\centering
\begin{tabularx}{\textwidth}{lXl}

       \hspace{6cm} &  & \hspace{6cm} \\
\cline{1-1}\cline{3-3}
Ort, Datum  &  & Unterschrift
\end{tabularx}
\end{table}								% Erkl�rung der Eigenarbeit
% \chapter*{Aufgabenstellung der Arbeit}
%\vspace*{4cm}

{\large \textbf{Thema:} \parbox[t]{0.8\textwidth}{Thema der Arbeit}}
\newline

Hier wird die Aufgabenstellung beschrieben. Die Notwendigkeit dieser hängt vom Betreuer ab.
 									% Aufgabenstellung (optional)
% \chapter*{Vorwort}

%\vspace*{4cm}

An dieser Stelle kann ein Vorwort eingefügt werden.									% Vorwort (optional)
\cleardoublepage
\begin{spacing}{1.15}									% evtl. kleinerer Zeilenabstand im IV, AV, TV
\pdfbookmark[1]{Inhaltsverzeichnis}{toc}					% Inhaltsverzeichnis bei den Lesezeichen rein
\tableofcontents 											% Inhaltsverzeichnis erzeugen
\listoffigures   											% Abbildungsverzeichnis
\listoftables   											% Tabellenverzeichnis
\end{spacing}

\addchap{Abstract}

Hierbei handelt es sich um einen, in englischer Sprache verfassten, kurzen Abriss der
Arbeit (Abstract). Der Umfang beträgt eine halbe bis eine Seite. Alternativ zu dem
Abstract kann eine in englischer Sprache verfasste Zusammenfassung (Summary) nach
der deutschen Zusammenfassung stehen. Im Hauptseminar kann dieser Teil entfallen.									% Abstract

\mainmatter												% Hauptteil beginnt
\chapter{Einleitung}
\label{kap:einleitung}

Dieses Kapitel dient zur Hinführung an das Thema. Hier ist herauszuarbeiten, warum das
Thema von Interesse und Bedeutung ist. Des Weiteren ist kurz der Aufbau der Arbeit
anzusprechen. Der Umfang des Kapitels sollte eine oder wenige Seiten betragen.									% Einleitung (Bezeichnung passend zur Arbeit!)
\chapter{Hauptteil}
\label{kap:hauptteil}

Hier ist das eigentliche Thema zu bearbeiten.

\section{Einbindung von Bildern}
\label{kap:einbindungbilder}
Abbildungen sind mit Hilfe des Pakets \textit{graphicx} einzufügen. Sie können im PDF-Format durch die Nutzung des folgenden Codes implementiert werden.
\begin{verbatim}
\begin{figure}[!htb]\centering
 \includegraphics*[width = \textwidth]{beispiel}
 \caption{Beispiel für die Einbindung eines Bildes}
 \label{abb:beispiel}
\end{figure}
\end{verbatim}
Das Ergebnis ist die Anzeige des Bildes, mittig, wie der Text breit mit der angegebenen Unterschrift. Alternativ kann bei der Breite eine absolute Angabe in mm erfolgen. Über das label \textit{abb:beispiel} kann das Bild referenziert werden.
\begin{figure}[!htb]\centering
 \includegraphics*[width = \textwidth]{beispiel}
 \caption{Beispiel für die Einbindung eines Bildes}
 \label{abb:beispiel}
\end{figure}
Um auf das Bild \ref{abb:beispiel} zu verweisen, bedient man sich der folgenden Funktion:
\begin{verbatim}
\ref{abb:beispiel}
\end{verbatim}
Die referenzierte Nummerierung erfolgt Kapitelweise. Will man weiterhin eine Quelle in der Bildunterschrift angeben, so ist darauf zu achten, dass die Einbindung der Bildunterschriften durch ein optionales Element (eingeschlossen in eckigen Klammern) erweitert wird, welches die Beschriftung für das Abbildungsverzeichnis enthält:
\begin{verbatim}
 \caption[Beispiel...]{Beispiel... , aus \cite{schwab}}
\end{verbatim}
Diese Variante verhindert, dass LaTex die Quellen bereits im Abbildungsverzeichnis zu zählen anfängt.

\section{Einbindung von Tabellen}
\label{kap:einbindungtabellen}
Die Tabellen sollen mit Hilfe des Pakets \textit{tabularx} eingebunden werden. Im Folgenden ist ein Beispiel für die Einbindung von Tabellen aufgeführt. Mit
\begin{verbatim}
\begin{table}[!htb]
\centering
\caption{Beispiel einer Tabelle}
\label{tab:tabelle1}
\begin{tabularx}{\textwidth}{|X|c|c|c|c|c|c|c|c|}
\hline
        & Spalte 1 & Spalte 2 & Spalte 3 & Spalte 4 & Spalte 5 \\
\hline
Zeile 1 &          &          &          &          &          \\
\hline
Zeile 2 &          &          &          &          &          \\
\hline
Zeile 3 &          &          &          &          &          \\
\hline
\end{tabularx}
\end{table}
\end{verbatim}
ergibt sich die folgende Tabellenausgabe \ref{tab:tabelle1}.
\begin{table}[!htb]
\centering
\caption{Beispiel einer Tabelle}
\label{tab:tabelle1}
\begin{tabularx}{\textwidth}{|X|c|c|c|c|c|c|c|c|}
\hline
 				& Spalte 1 	& Spalte 2 	& Spalte 3 	& Spalte 4 	& Spalte 5 \\
\hline
Zeile 1 & 					& 					& 					& 					& \\
\hline
Zeile 2 & 					& 					& 					& 					& \\
\hline
Zeile 3 & 					& 					& 					& 					& \\
\hline
\end{tabularx}
\end{table}
Bei Angabe von Quellen in der Tabellenüberschrift ist ähnlich wie im Kapitel \ref{kap:einbindungbilder} zu verfahren.

\section{Eingabe von Gleichungen}
\label{kap:einbindunggleichungen}


\section{Quellenangabe}
\label{kap:quellenangabe}
Die Angabe von Quellen ist mit folgendem Code als Beispiel möglich:
\begin{verbatim}
\cite{schwab}
\end{verbatim}
Die Ausgabe wird an der Stelle eingefügt, an der man es einsetzt. Zum Beispiel hier: \cite{schwab}. Die Zahlen werden nach dem Vorkommen im Text vergeben und durchnummeriert. D.h. die nächste Angabe erhält die Nummer zwei und zwar hier: \cite{oedingoswald}. Die Einträge in der Literaturdatenbank können manuell in der Datei \textit{literatur.bib} oder mit der Software \textit{JabRef} angepasst werden.									% Hauptteil (Bezeichnung passend zur Arbeit!)
\chapter{Zusammenfassung}
\label{kap:zusammenfassung}

In der Zusammenfassung werden die Ergebnisse der Arbeit kurz zusammengefasst. Der
Umfang beträgt ca. eine Seite.							% Zusammenfassung (Bezeichnung passend zur Arbeit!)

\appendix												% Anhang
\chapter{Anhang Teil 1}
%\chapter{Anhang Teil 2}
%\chapter{Anhang Teil 3}

									% Anhangsinhalt
\addchap{Symbol- und Abkürzungsverzeichnis}
\label{kap:symbole}

Sollten in einer Ausarbeitung viele Abkürzungen und Formelzeichen auftreten, so
empfiehlt es sich, diese gesondert in einem Kapitel aufzufuhren. Dieses kann auch nach
dem Inhaltsverzeichnis (Abbildungs- und Tabellenverzeichnis) folgen.
\begin{tabbing}
XXXXXXXXXXXXX \= XXXX \= XXXXXXXXXXXXXXXXXXXXXXXXXXXXXXXXXXXXXXXXXXXXXXXXX \kill
A									\>										\> Abkürzung für ... \\
$\mathrm{cos(\varphi)}$				\>										\> Leistungsfaktor \\
$U_\mathrm{S}$						\> V									\> Betrag der Statorspannung \\
$\varphi$							\> rad									\> Winkel zwischen Spannung und Strom \\
\end{tabbing}									% Symbol- und Abk�rzungsverzeichnis

\begin{spacing}{1.5}          								% evtl. kleinerer Zeilenabstand im LV
%\nocite{*}                              							% alle Literaturquellen einbinden, sonst werden nur die zitierten
                                           								% Quellen im Literaturverzeichnis angezeigt.
\bibliographystyle{unsrtdin}		         					% Darstellung numerisch nach Vorkommen im Text
\bibliography{literatur}
\end{spacing}

\end{document}
%------------------------- Ende des Dokuments -----------------------